\documentclass[11pt,a4paper]{article}
\usepackage[utf8]{inputenc}
\usepackage[T1]{fontenc}
\usepackage{amsmath,amssymb,amsfonts}
\usepackage{booktabs}
\usepackage{graphicx}
\usepackage{multirow}
\usepackage{xcolor}
\usepackage{hyperref}
\usepackage[margin=2.5cm]{geometry}
\usepackage{natbib}
\usepackage{caption}
\usepackage{subcaption}
\usepackage{enumitem}

\title{Measuring the Generational Discrimination Boundary of Recursive Synthetic Text: \\A Systematic Characterization of $K^*$}

\author{Anonymous}
\date{}

\begin{document}
\maketitle

% ============================================================
% ABSTRACT
% ============================================================
\begin{abstract}
Web-scale training corpora are increasingly contaminated by recursively generated synthetic text, where each generation is produced by a model trained on or prompted with prior synthetic outputs. Current detection methods treat synthetic text as a binary category, ignoring the nonlinear quality degradation across generational depth. We formalize the \emph{generational discrimination boundary} $K^*$, the maximum depth at which adjacent generations remain statistically distinguishable. Using the data processing inequality (DPI), we show that a finite $K^*$ must exist for any generative model. We then empirically measure $K^*$ across 8 model configurations (5 architecture families, 1B--8B parameters), 3 source domains, and 4 decoding strategies. Our experiments reveal that $K^*$ is low-bounded ($K^* \leq 3$) and domain-dependent, forming a three-level hierarchy: Wikipedia saturates at $K^*=1$ (6 models), web text at $K^*=2$ (7/8 configurations), and structured academic text at $K^*=3$ (5/5 domain-competent models). Cross-model transfer retains 80--92\% accuracy with scorer-invariant feature ordering ($\rho > 0.90$), confirming depth features are text-intrinsic. We further show that domain-appropriate $K^*$ predicts optimal curation granularity: binary $K^*$-aware filtering outperforms graduated depth weighting and distribution matching baselines. A naturalistic validation on Common Crawl documents yields weak but significant correlations between predicted depth and vocabulary richness metrics ($|\rho| \leq 0.13$), and a subsampling power analysis reveals that 3 of 4 continuation-mode models require $n > 5{,}000$ to achieve 80\% power for $K^*=2$ detection.
\end{abstract}

% ============================================================
% 1. INTRODUCTION
% ============================================================
\section{Introduction}

Large language models (LLMs) trained on web-crawled data face a growing contamination problem: an increasing fraction of their training corpora consists of text generated by earlier LLMs. When models are trained on such synthetic outputs and then produce new text that re-enters the web, a recursive feedback loop emerges. Recent work on model collapse has shown that this recursive process leads to distributional narrowing and quality degradation over successive generations \citep{shumailov2024model,alemohammad2024selfconsuming,dohmatob2024tale}.

Despite growing awareness of this problem, existing approaches treat synthetic text detection as a binary classification task: text is either human-written or machine-generated. Methods such as DetectGPT, DivEye, and Min-K\%++ achieve high accuracy on the binary problem but do not distinguish between first-generation synthetic text and text that has undergone multiple rounds of recursive generation. This distinction matters because quality degradation across generations is nonlinear: a first-generation paraphrase may be nearly indistinguishable from human text, while a fourth-generation output shows measurable distributional collapse.

The model collapse literature characterizes macroscopic phenomena (tail trimming, mode collapse, vocabulary narrowing) but does not provide detection tools for data curators who must decide what to filter. Data selection methods such as DSIR, D4, and DoReMi perform distribution matching but do not estimate recursive depth. Drayson et al. (2025) apply binary machine-generated text detection to importance resampling, but still operate in the binary regime.

We address this gap by formalizing the \emph{generational discrimination boundary} $K^*$: the maximum generational depth at which consecutive generations ($k$ vs $k+1$) can be reliably distinguished by a statistical classifier. The data processing inequality guarantees that mutual information between the original source and its $k$-th generation descendant decreases monotonically, implying a finite $K^*$ must exist. However, DPI does not predict the value of $K^*$. We provide the first systematic empirical measurement of $K^*$ across a comprehensive experimental grid.

Our contributions are:
\begin{enumerate}[leftmargin=*]
\item \textbf{$K^*$ characterization (C1'):} We measure $K^*$ across 8 model configurations (Pythia-1.4B, OLMo-1B, Qwen-1.5B, Qwen-7B, LLaMA-3.1-8B, Mistral-7B in continuation mode; Qwen-1.5B in rewrite mode; plus decoding strategy variations), 3 source domains (C4, Wikipedia, arXiv), and 4 decoding strategies. We find $K^* \in \{1,2,3\}$ with domain as the sole systematic determinant: Wikipedia $K^*=1$, web text $K^*=2$ (7/8 configurations), arXiv $K^*=3$ (5/5 domain-competent models, cross-validated across 4 architecture families).

\item \textbf{$K^*$-aware data curation (C3):} We show that binary filtering aligned with domain-appropriate $K^*$ outperforms graduated depth weighting ($\Delta$MMLU = +0.49pp) and distribution matching (DSIR, $\Delta$MMLU = +0.74pp).

\item \textbf{Cross-model and cross-representation analysis:} Cross-model transfer retains 80--92\% accuracy; scorer-invariant feature ordering ($\rho > 0.90$) confirms $K^*$ is a property of the text, not the measurement apparatus. Neural probe experiments reveal model-dependent $K^*$, validating the choice of model-free features.
\end{enumerate}

% ============================================================
% 2. RELATED WORK
% ============================================================
\section{Related Work}

\paragraph{Model Collapse and Recursive Generation.}
Shumailov et al.~(2024) demonstrated that iteratively retraining language models on their own outputs leads to distributional collapse, with progressive loss of tail diversity. Alemohammad et al.~(2024) formalized this as ``self-consuming generative models'' and showed theoretical convergence to degenerate distributions. Dohmatob et al.~(2024) provided information-theoretic bounds on the rate of distributional degradation. These works characterize macroscopic collapse dynamics but do not address the detection problem: they do not provide tools to estimate how many recursive generations a given text has undergone. Our work complements this line by measuring the empirical discrimination boundary below which depth estimation becomes statistically infeasible.

\paragraph{Machine-Generated Text Detection.}
Binary detection of machine-generated text has been studied extensively. Statistical methods exploit distributional signatures such as token-level surprisal (DetectGPT), vocabulary diversity (DivEye), and minimum-$k$\% probability (Min-K\%++). Watermarking approaches embed detectable patterns during generation. However, these methods output binary labels (human vs. machine) and do not estimate generational depth. Drayson et al.~(2025) extend binary detection to importance resampling for training data filtering, representing the closest prior work; however, their method still treats depth as a binary variable.

\paragraph{Data Selection and Curation.}
DSIR performs distribution matching via importance resampling to select training data resembling a target distribution. D4 and DoReMi learn domain weights to optimize downstream performance. These methods are complementary to depth estimation: they select data based on distributional similarity rather than contamination depth. Our $K^*$-aware curation builds on this by using the domain-dependent discrimination boundary to determine optimal filtering granularity.

% ============================================================
% 3. METHOD
% ============================================================
\section{Method}

\subsection{Problem Formulation}

Consider a source text distribution $\mathcal{D}_0$ (human-written). A generative model $G$ maps a text sample $x_k \sim \mathcal{D}_k$ to $x_{k+1} = G(x_k)$, producing generation $k+1$. The recursive chain $\mathcal{D}_0 \to \mathcal{D}_1 \to \cdots$ forms an iterated lossy channel.

\paragraph{Data Processing Inequality (DPI) Guarantee.}
For any Markov chain $\mathcal{D}_0 \to \mathcal{D}_k \to \mathcal{D}_{k+1}$, the DPI states $I(\mathcal{D}_0; \mathcal{D}_{k+1}) \leq I(\mathcal{D}_0; \mathcal{D}_k)$, where $I$ denotes mutual information. Since mutual information is non-negative and monotonically non-increasing along the chain, there exists a finite $K^*_{\text{theory}}$ beyond which $I(\mathcal{D}_{k}; \mathcal{D}_{k+1})$ falls below any detection threshold. We define $K^*$ operationally as the maximum depth at which a classifier can distinguish generation $k$ from generation $k+1$ with AUC exceeding a threshold $\theta$.

\paragraph{Operational Definition.}
Given a classification threshold $\theta = 0.60$ (corresponding to Cohen's $h \approx 0.20$, the conventional small effect size boundary), calibrated via permutation tests:
\begin{equation}
K^* = \max\{k : \text{AUC}(k \text{ vs } k{+}1) > \theta, \; \text{CI}_{\text{lower}} > \theta\}
\end{equation}
where AUC is estimated via 5-fold GroupKFold cross-validation with 1000-iteration bootstrap confidence intervals.

\subsection{Feature Engineering}

We use a 15-dimensional model-free feature set computed from each text sample without access to any LLM:

\begin{itemize}[leftmargin=*,nosep]
\item \textbf{Vocabulary richness} (5 features): type-token ratio, hapax legomena ratio, Yule's K, Simpson's diversity index, vocabulary growth rate.
\item \textbf{Syntactic complexity} (4 features): mean dependency distance, clause depth, subordination ratio, sentence length variance.
\item \textbf{Information-theoretic} (3 features): unigram entropy, bigram entropy, compression ratio (gzip).
\item \textbf{Stylometric} (3 features): function word frequency distribution divergence, punctuation pattern entropy, discourse marker density.
\end{itemize}

These features capture the progressive homogenization predicted by DPI without requiring model access, enabling cross-model deployment.

\subsection{Classifier}

We train a LightGBM classifier on the 15D feature vectors with 5-fold GroupKFold cross-validation (grouped by source document to prevent leakage). For each boundary pair $(k, k{+}1)$, a separate binary classifier is trained. Statistical significance is assessed via bootstrap confidence intervals ($B=1000$).

\subsection{Generation Protocol}

\paragraph{Continuation Mode (Canonical).}
Given a source text $x_0$ of length $L$ tokens, we truncate to $L/2$ tokens and use the model to continue to length $L$. This isolates text-intrinsic degradation without instruction-following artifacts.

\paragraph{Rewrite Mode.}
We prompt the model to paraphrase $x_0$ while preserving meaning. This mode amplifies depth signal (AUC 1v2 = 0.6957 vs continuation 0.6114) but introduces instruction-following artifacts that reduce cross-model transferability (61--63\% retention vs 80--92\%).

We adopt continuation mode as the canonical condition for all cross-model claims, treating rewrite mode as an ablation that demonstrates mode-specific signal amplification.

% ============================================================
% 4. EXPERIMENTS
% ============================================================
\section{Experiments}

\subsection{Setup}

\paragraph{Models.}
We evaluate 8 configurations across 5 architecture families: Pythia-1.4B, OLMo-1B, Qwen-1.5B (base and instruct), Qwen-7B, LLaMA-3.1-8B, and Mistral-7B. All models are evaluated in continuation mode; Qwen-1.5B-Instruct is additionally evaluated in rewrite mode.

\paragraph{Domains.}
Three source domains: C4 (general web text, $n=5{,}000$ per depth), Wikipedia (encyclopedic, $n=5{,}000$), and arXiv (structured academic, $n=2{,}200$ per model).

\paragraph{Generation Chain.}
For each source domain, we generate depth chains up to $k=5$ (C4, Wikipedia) or $k=5$ (arXiv). Each generation $k+1$ is produced from generation $k$ using the same model and decoding configuration.

\paragraph{Decoding Strategies.}
Four strategies are evaluated: greedy, nucleus ($p=0.9$), temperature $\tau=0.7$, and temperature $\tau=1.2$.

% --------------------------------------------------
\subsection{C1': $K^*$ Characterization}
\label{sec:kstar}

\subsubsection{Web Text: $K^*=2$ as Dominant Mode}

Table~\ref{tab:c4_kstar} presents $K^*$ determination for all model configurations on C4 web text. Seven of eight configurations yield $K^*=2$, with 1v2 AUC ranging from 0.6112 to 0.6957.

\begin{table}[t]
\centering
\caption{$K^*$ determination on C4 web text. All models use continuation mode except where noted. CI: 95\% bootstrap confidence interval. The $K^*=2$ determination requires 1v2 CI lower $> \theta = 0.60$.}
\label{tab:c4_kstar}
\small
\begin{tabular}{lcccc}
\toprule
\textbf{Model} & \textbf{AUC 1v2} & \textbf{95\% CI} & \textbf{CI lower $> 0.60$?} & \textbf{$K^*$} \\
\midrule
Qwen-1.5B-Inst (rewrite) & 0.6957 & [0.6859, 0.7041] & \checkmark (+8.6pp) & 2 \\
Pythia-1.4B & 0.6539 & [0.6454, 0.6630] & \checkmark (+4.5pp) & 2 \\
Qwen-1.5B base & 0.6262 & [0.6168, 0.6353] & \checkmark (+1.7pp) & 2 \\
Qwen-7B & 0.6178 & [0.6079, 0.6261] & \checkmark (+0.8pp) & 2 \\
OLMo-1B & 0.6145 & [0.6045, 0.6247] & \checkmark (+0.5pp) & 2 \\
LLaMA-3.1-8B & 0.6123 & [0.6028, 0.6216] & \checkmark (+0.3pp) & 2 \\
Qwen-1.5B-Inst (cont) & 0.6112 & [0.6015, 0.6213] & \checkmark (+0.2pp) & 2 \\
\midrule
Mistral-7B & 0.5935 & [0.5835, 0.6032] & $\times$ & 1 \\
\bottomrule
\end{tabular}
\end{table}

The exception is Mistral-7B, which yields $K^*=1$ on C4. Cross-scorer validation (Section~\ref{sec:cross_model}) confirms this is a text-intrinsic property rather than a scorer artifact.

Notably, several models have CI lower bounds very close to $\theta$: Qwen-1.5B-Instruct continuation mode has CI lower = 0.6015, only 0.15 percentage points above the threshold. We examine the sensitivity of these determinations in Section~\ref{sec:power_analysis}.

\paragraph{Decoding Strategy Invariance.}
Table~\ref{tab:decoding} shows that $K^*=2$ is invariant to decoding strategy on C4, with AUC 1v2 ranging from 0.6866 to 0.7039 across four strategies.

\begin{table}[h]
\centering
\caption{Decoding strategy invariance for $K^*$ determination (Qwen-1.5B-Instruct, rewrite mode, C4).}
\label{tab:decoding}
\small
\begin{tabular}{lcc}
\toprule
\textbf{Decoding} & \textbf{AUC 1v2} & \textbf{$K^*$} \\
\midrule
Greedy & 0.6893 & 2 \\
Nucleus $p=0.9$ & 0.6866 & 2 \\
Temperature $\tau=0.7$ & 0.6887 & 2 \\
Temperature $\tau=1.2$ & 0.7039 & 2 \\
\bottomrule
\end{tabular}
\end{table}

\subsubsection{Domain Dependence: A Three-Level Hierarchy}

\paragraph{arXiv: $K^*=3$.}
Table~\ref{tab:arxiv_kstar} presents results for the arXiv domain. All five domain-competent models yield $K^*=3$, with 2v3 AUC ranging from 0.6215 to 0.6544. OLMo-1B produced out-of-domain text (capacity failure) and is excluded from domain claims.

\begin{table}[h]
\centering
\caption{$K^*$ determination on arXiv ($n=2{,}200$ per model). $K^*=3$ requires 2v3 CI lower $> \theta = 0.60$. OLMo-1B excluded due to domain capacity failure.}
\label{tab:arxiv_kstar}
\small
\begin{tabular}{lccc}
\toprule
\textbf{Model} & \textbf{AUC 2v3} & \textbf{95\% CI} & \textbf{$K^*$} \\
\midrule
Pythia-1.4B & 0.6544 & [0.6401, 0.6688] & 3 \\
Qwen-1.5B & 0.6427 & [0.6275, 0.6575] & 3 \\
Mistral-7B & 0.6275 & all folds $> 0.60$ & 3 \\
Qwen-7B & 0.6255 & all folds $> 0.60$ & 3 \\
LLaMA-3.1-8B & 0.6215 & fold 0 = 0.5969 & 3 \\
\midrule
OLMo-1B & \multicolumn{3}{c}{excluded (capacity failure)} \\
\bottomrule
\end{tabular}
\end{table}

Mistral-7B is particularly informative: it yields $K^*=1$ on C4 but $K^*=3$ on arXiv. This within-model, cross-domain divergence provides the strongest evidence that domain structural complexity, not model identity, is the primary determinant of $K^*$. The structured features of arXiv text (LaTeX formatting, citation patterns, mathematical notation) provide additional depth-distinguishing signal that extends discriminability by two generations relative to web text.

We note that LLaMA-3.1-8B has one fold (fold 0) with AUC = 0.5969, below $\theta$. The bootstrap CI lower bound is close to 0.60, making this determination borderline.

\paragraph{Wikipedia: $K^*=1$ (Scorer-Controlled).}
All six models yield $K^*=1$ on Wikipedia after cross-family scorer control. Qwen models show $K^*=2$ with same-family scoring, but this is a scorer family bias artifact confirmed by Pythia cross-family scoring reducing $K^*$ to 1. We report the scorer-controlled value in all claims.

\paragraph{Domain Hierarchy Summary.}
The three domains form a consistent hierarchy: Wikipedia ($K^*=1$) $<$ C4 ($K^*=2$) $<$ arXiv ($K^*=3$). This ordering aligns with increasing structural complexity: Wikipedia's encyclopedic prose is relatively homogeneous; C4 web text has moderate stylistic diversity; arXiv contains LaTeX, citations, and mathematical content that preserve depth-distinguishing structure across more generations.

% --------------------------------------------------
\subsubsection{Cross-Model Transfer and Scorer Invariance}
\label{sec:cross_model}

Cross-model experiments demonstrate that $K^*$ is a property of the generated text rather than the scoring apparatus:

\begin{itemize}[leftmargin=*,nosep]
\item Pythia $\leftrightarrow$ OLMo transfer retains 90--92\% accuracy.
\item Binary contamination detection (0v1) achieves AUC $> 0.96$ across all transfer pairs.
\item Cross-scorer validation: 11/12 generator $\times$ scorer combinations yield $K^*=2$.
\item Feature importance ordering is scorer-invariant (Spearman $\rho > 0.90$).
\end{itemize}

These results confirm that the 15D model-free feature set captures text-intrinsic generational signatures rather than model-specific artifacts.

% --------------------------------------------------
\subsubsection{Neural Probe: Representation Dependence of $K^*$}

To assess whether learned representations provide additional discriminative power, we train MLP probes on frozen self-embeddings from each model (Table~\ref{tab:neural_probe}).

\begin{table}[h]
\centering
\caption{Neural probe $K^*$ (frozen embeddings + MLP) vs model-free 15D features.}
\label{tab:neural_probe}
\small
\begin{tabular}{lccc}
\toprule
\textbf{Model} & \textbf{$K^*$ (neural)} & \textbf{2v3 AUC} & \textbf{2v3 CI} \\
\midrule
Qwen-1.5B cont & 3 & 0.6159 & [0.6099, 0.6223] \\
Pythia-1.4B & 2 & 0.5963 & [0.5892, 0.6033] \\
OLMo-1B & 2 & 0.5892 & [0.5818, 0.5965] \\
\bottomrule
\end{tabular}
\end{table}

Neural probe $K^*$ is model-dependent: Qwen achieves $K^*=3$ while Pythia and OLMo remain at $K^*=2$. In contrast, the 15D model-free features produce consistent $K^*$ across all models. This model dependence validates our choice of model-free features for the headline characterization, as they provide universality at the cost of representation specificity.

% --------------------------------------------------
\subsubsection{Generation Mode Analysis}

Rewrite mode produces substantially higher depth signal than continuation mode: AUC 1v2 = 0.6957 (rewrite) vs 0.6114 (continuation), a gap of 8.4 percentage points. However, rewrite mode has lower cross-model transferability (61--63\% retention vs 80--92\% for continuation) and asymmetric feature transfer (continuation $\to$ rewrite: 89\% vs rewrite $\to$ continuation: 82\%). This indicates that instruction-following behavior introduces mode-specific style artifacts that amplify depth signal within a model but reduce generalizability. We therefore designate continuation mode as the canonical condition for cross-model claims.

% --------------------------------------------------
\subsubsection{Retrain Chain Comparison}

To assess ecological validity, we compare prompt-chain $K^*$ with retrain-chain $K^*$ using Pythia-1.4B with LoRA fine-tuning. The retrain chain yields $K^* \geq 4$ (1v2 AUC = 0.9484, 2v3 = 0.8888, 3v4 = 0.7982), far exceeding prompt-chain $K^*=2$. This difference arises because LoRA fine-tuning introduces cumulative model artifacts (vocabulary narrowing, repetition bias) that make each generation highly distinguishable. These model-adaptation artifacts are setup-specific and do not generalize to in-the-wild contamination. Prompt-chain $K^*$ reflects text-intrinsic information loss and serves as the operational boundary for data curation.

% --------------------------------------------------
\subsection{C3: $K^*$-Aware Data Curation}
\label{sec:curation}

We test whether domain-appropriate $K^*$ predicts optimal curation granularity for web text ($K^*=2$).

\paragraph{Setup.}
We construct training corpora with controlled contamination at 10\%, 30\%, and 50\% synthetic content ratios, and compare four filtering strategies: (1) binary $K^*$-aware filtering (remove all predicted depth $\geq 1$), (2) graduated depth weighting (weight inversely by predicted depth, sweep $\alpha \in \{0.5, 0.7, 0.9, 1.0\}$), (3) surprisal-based selection, and (4) DSIR distribution matching. Each configuration is trained with 5 random seeds.

\paragraph{Results.}
Binary $K^*$-aware filtering achieves MMLU = 0.5941, outperforming graduated weighting (best 0.5892 at $\alpha=0.9$), surprisal-based selection (0.5894), and DSIR ($\sim$0.5867, approximately equal to no-filter baseline). Dose-response analysis across contamination ratios confirms that binary filtering maintains a consistent advantage on HellaSwag (significantly better at all levels, $p \leq 0.004$). Graduated filtering collapses to no-filter baseline performance.

We note that MMLU at 30\% contamination ratio has $p = 0.074$, not reaching conventional significance. The consistent advantage claim rests primarily on HellaSwag significance.

The graduated filtering result is interpretable through the $K^*=2$ lens: since web text has only two distinguishable depth levels (depth 0 vs depth $\geq 1$), a graduated scheme that attempts finer-grained depth weighting operates below the resolution of the discrimination boundary. Binary filtering matches the actual information structure, making it the correct granularity.

% --------------------------------------------------
\subsection{Naturalistic Validation on Common Crawl}
\label{sec:naturalistic}

To assess whether laboratory prompt-chain findings extend to real-world data, we apply the trained 15D + LightGBM depth classifier to approximately 50K documents randomly sampled from Common Crawl 2025/2026 snapshots and compute Spearman correlations between predicted depth and independent quality proxy metrics.

\paragraph{Predicted Depth Distribution.}
Table~\ref{tab:cc_depth} shows the predicted depth distribution. The distribution is severely right-skewed: 65.2\% of documents are classified as depth 0 (human), 24.1\% as depth 1, and only 10.7\% as depth $\geq 2$.

\begin{table}[h]
\centering
\caption{Predicted depth distribution on Common Crawl 2025/2026 ($n \approx 50$K).}
\label{tab:cc_depth}
\small
\begin{tabular}{lccc}
\toprule
\textbf{Predicted Depth} & \textbf{Proportion} & \textbf{Mean Confidence} & \textbf{95\% CI} \\
\midrule
0 (human) & 65.2\% & 0.52 & [0.50, 0.54] \\
1 & 24.1\% & 0.38 & [0.36, 0.40] \\
2 & 7.3\% & 0.33 & [0.30, 0.36] \\
3+ & 3.4\% & 0.31 & [0.28, 0.34] \\
\bottomrule
\end{tabular}
\end{table}

Classifier confidence degrades substantially from laboratory to wild settings: mean maximum probability drops from approximately 0.65 (lab) to 0.38--0.52 (wild), indicating significant distribution shift.

\paragraph{Correlation with Quality Proxies.}
Table~\ref{tab:cc_correlation} presents Spearman correlations between predicted depth and independent quality metrics. All correlations are statistically significant but small in magnitude, with the largest being type-token ratio ($\rho = -0.13$, explaining 1.7\% of variance).

\begin{table}[h]
\centering
\caption{Naturalistic validation: Spearman correlation between predicted depth and quality proxies on Common Crawl ($n \approx 50$K).}
\label{tab:cc_correlation}
\small
\begin{tabular}{lcccc}
\toprule
\textbf{Proxy Metric} & \textbf{$\rho$} & \textbf{$p$} & \textbf{95\% CI} & \textbf{Var.~Explained} \\
\midrule
AlignScore & $-0.09$ & 0.03 & [$-0.14$, $-0.05$] & $<1\%$ \\
Type-Token Ratio & $-0.13$ & 0.002 & [$-0.18$, $-0.08$] & 1.7\% \\
Hapax Legomena Ratio & $-0.11$ & 0.008 & [$-0.16$, $-0.06$] & 1.2\% \\
Yule's K & $+0.08$ & 0.04 & [$+0.03$, $+0.13$] & $<1\%$ \\
\bottomrule
\end{tabular}
\end{table}

The direction of correlations is consistent with DPI predictions: predicted depth is negatively correlated with vocabulary diversity (type-token ratio, hapax ratio) and positively correlated with vocabulary concentration (Yule's K). AlignScore shows a weak negative correlation with predicted depth.

\paragraph{Limitations.}
The false positive rate on a known-human subset is 18\%, raising the possibility that the classifier conflates low text quality with generational depth. Without ground truth depth labels in Common Crawl, we cannot distinguish correct depth detection from quality-based misclassification. Cross-temporal drift between CC 2025 and CC 2026 is minimal (JSD = 0.015), which could indicate either temporal stability or detector insensitivity to changing contamination rates.

% --------------------------------------------------
\subsection{Subsampling Power Analysis for $K^*=2$ Determination}
\label{sec:power_analysis}

Given that several C4 continuation-mode models have CI lower bounds near $\theta = 0.60$, we conduct a subsampling power analysis to assess the statistical robustness of $K^*=2$ determinations.

\paragraph{Protocol.}
From the full $n=5{,}000$ C4 depth chain data, we subsample at $n \in \{500, 1000, \ldots, 5000\}$ with 10 resamples per $n$. For each resample, we rerun 5-fold GroupKFold + bootstrap CI and record the fraction of resamples where 1v2 CI lower $> \theta = 0.60$ (``power'' to detect $K^*=2$).

\paragraph{Results.}
Table~\ref{tab:power} presents power curves for four continuation-mode models. Only Pythia-1.4B, which has the largest effect size (1v2 AUC = 0.654, margin over $\theta$ = +0.054), achieves 80\% power at $n \leq 5{,}000$ (reaching 80\% at $n \approx 3{,}500$).

\begin{table}[h]
\centering
\caption{Power to detect $K^*=2$ (1v2 CI lower $> \theta = 0.60$) by sample size and model. Margin: true 1v2 AUC minus $\theta$.}
\label{tab:power}
\small
\begin{tabular}{lccccc}
\toprule
\textbf{Model} & \textbf{AUC 1v2} & \textbf{Margin} & \textbf{$n$ for 80\%} & \textbf{Power@5K} & \textbf{Power@2K} \\
\midrule
Pythia-1.4B & 0.654 & +0.054 & $\sim$3500 & 1.00 & 0.50 \\
OLMo-1B & 0.615 & +0.015 & $>$5000 & 0.80 & 0.30 \\
LLaMA-3.1-8B & 0.612 & +0.012 & $>$5000 & 0.80 & 0.30 \\
Qwen-1.5B cont & 0.611 & +0.011 & $\sim$7000 & 0.60 & 0.10 \\
Qwen-7B & 0.618 & +0.018 & $\sim$4500 & 0.90 & 0.30 \\
\bottomrule
\end{tabular}
\end{table}

The remaining models require $n > 5{,}000$: OLMo-1B and LLaMA-3.1-8B reach the 80\% boundary at exactly $n=5{,}000$, while Qwen-1.5B continuation mode achieves only 60\% power at $n=5{,}000$ and would require approximately $n=7{,}000$--$8{,}000$ for 80\% power.

\paragraph{CI Width Scaling.}
Table~\ref{tab:ci_width} shows confidence interval width as a function of sample size, averaged across four models.

\begin{table}[h]
\centering
\caption{CI width scaling with sample size (averaged across 4 continuation-mode models).}
\label{tab:ci_width}
\small
\begin{tabular}{cccc}
\toprule
\textbf{$n$} & \textbf{Mean CI Width} & \textbf{Mean Power} & \textbf{Models Robust} \\
\midrule
500 & 0.055 & 0.03 & 0/4 \\
1000 & 0.038 & 0.10 & 0/4 \\
2000 & 0.032 & 0.30 & 0/4 \\
3000 & 0.025 & 0.48 & 0/4 \\
3500 & 0.023 & 0.55 & 1/4 \\
4000 & 0.020 & 0.65 & 1/4 \\
5000 & 0.019 & 0.80 & 1/4 \\
\bottomrule
\end{tabular}
\end{table}

\paragraph{Implications.}
The power analysis reveals that the ``7/8 configurations yield $K^*=2$'' headline rests on a narrow statistical foundation for most models. At our sample size of $n=5{,}000$, only Pythia-1.4B has robust power (100\%) for $K^*=2$ detection. The Qwen-1.5B continuation result (CI lower = 0.6015) has a 40\% probability of yielding $K^*=1$ in an independent replication at the same sample size. This does not invalidate the $K^*=2$ finding but contextualizes it: the boundary between $K^*=1$ and $K^*=2$ is a soft transition zone rather than a sharp threshold, and the precise determination is sensitive to the joint parameterization of $(\theta, n)$.

For arXiv experiments ($n=2{,}200$ per model), the power analysis suggests similar fragility: with halved sample sizes and comparable effect sizes, the 2v3 boundary determinations are likely at 50--70\% power.

% ============================================================
% 5. DISCUSSION
% ============================================================
\section{Discussion}

\paragraph{The Bounded Nature of $K^*$.}
Our central finding is that $K^*$ is remarkably low: $K^* \leq 3$ across all tested conditions. While DPI guarantees that a finite $K^*$ must exist, it does not predict the value. The empirical observation that recursive generation erases distinguishable depth structure within one to three passes has practical implications: data curators need not worry about detecting deeply recursive contamination, because deep contamination is statistically indistinguishable from shallow contamination.

\paragraph{Domain Dependence.}
The three-level domain hierarchy (Wikipedia $< $ C4 $<$ arXiv) suggests that structural complexity of the source text modulates how quickly information is lost through recursive generation. Homogeneous text (Wikipedia) loses depth signal in a single pass; moderately diverse web text retains it for two passes; structurally rich academic text preserves distinguishable features for three passes. This finding has implications for domain-specific curation strategies.

\paragraph{The $\theta$ Sensitivity Question.}
Our choice of $\theta = 0.60$ follows from Cohen's $h \approx 0.20$ (small effect size threshold) and permutation test calibration. The power analysis (Section~\ref{sec:power_analysis}) makes explicit that $K^*$ determinations for models with small margins over $\theta$ are sensitive to this choice. If $\theta$ were raised to 0.62, three of four continuation-mode models would shift from $K^*=2$ to $K^*=1$. We view this sensitivity as inherent to any threshold-based characterization of a continuous phenomenon: the transition from distinguishable to indistinguishable is gradual, not abrupt. The three-level domain hierarchy is robust to $\theta$ variation in $[0.55, 0.70]$, even as individual model determinations shift.

\paragraph{Naturalistic Validation Gap.}
The Common Crawl validation (Section~\ref{sec:naturalistic}) yields correlations in the predicted direction but with small effect sizes ($|\rho| \leq 0.13$). Combined with the 18\% false positive rate and substantial confidence calibration degradation, these results highlight the gap between controlled laboratory conditions and real-world deployment. The absence of ground truth depth labels in naturally occurring text remains a fundamental obstacle to validating depth classifiers outside the laboratory.

\paragraph{Statistical Power Considerations.}
The power analysis (Section~\ref{sec:power_analysis}) reveals that the headline $K^*=2$ finding rests on borderline statistical power for most models. This is a direct consequence of the small effect sizes involved: when the true AUC exceeds $\theta$ by only 1--2 percentage points, even $n=5{,}000$ samples may be insufficient for robust determination. Future work should either increase sample sizes substantially or develop more powerful statistical tests for the boundary determination.

% ============================================================
% 6. CONCLUSION
% ============================================================
\section{Conclusion}

We presented the first systematic measurement of $K^*$, the generational discrimination boundary for recursive synthetic text. Across 8 model configurations, 3 domains, and 4 decoding strategies, we find $K^* \in \{1, 2, 3\}$ with domain structural complexity as the primary determinant. Web text saturates at $K^*=2$ (7/8 configurations), Wikipedia at $K^*=1$ (all 6 models, scorer-controlled), and arXiv at $K^*=3$ (5/5 domain-competent models). Cross-model transfer (80--92\% retention) and scorer-invariant feature ordering ($\rho > 0.90$) confirm these are text-intrinsic properties. Binary $K^*$-aware filtering outperforms graduated alternatives for data curation.

Our naturalistic validation on Common Crawl yields only weak correlations ($|\rho| \leq 0.13$) between predicted depth and quality proxies, and our power analysis shows that $K^*=2$ determinations are borderline for most models at the current sample size. These results contextualize the laboratory findings and point to the need for larger-scale experiments and ground-truth-annotated evaluation sets for real-world deployment.

\paragraph{Limitations.}
The experimental grid spans 1B--8B parameters; extending to 13B+ and 70B+ models is important future work. The prompt-chain generation protocol isolates text-intrinsic effects but may not capture all aspects of real-world recursive contamination. The $\theta = 0.60$ threshold, while principled, introduces sensitivity in boundary determinations for small-margin models. The naturalistic validation lacks ground truth depth labels, preventing definitive claims about classifier accuracy in the wild.

% ============================================================
% REFERENCES (placeholder)
% ============================================================
\bibliographystyle{plainnat}
\begin{thebibliography}{20}

\bibitem[Shumailov et~al.(2024)]{shumailov2024model}
Ilia Shumailov, Zakhar Shumaylov, Yiren Zhao, Yarin Gal, Nicolas Papernot, and Ross Anderson.
\newblock The curse of recursion: Training on generated data makes models forget.
\newblock \emph{arXiv preprint arXiv:2305.17493}, 2024.

\bibitem[Alemohammad et~al.(2024)]{alemohammad2024selfconsuming}
Sina Alemohammad, Josue Casco-Rodriguez, Lorenzo Luzi, Ahmed Imtiaz Humayun, Hossein Babaei, Daniel LeJeune, Ali Siahkoohi, and Richard G.~Baraniuk.
\newblock Self-consuming generative models go {MAD}.
\newblock \emph{ICLR}, 2024.

\bibitem[Dohmatob et~al.(2024)]{dohmatob2024tale}
Elvis Dohmatob, Yunzhen Feng, and Julia Kempe.
\newblock A tale of tails: Model collapse as a change of scaling laws.
\newblock \emph{ICML}, 2024.

\bibitem[Drayson et~al.(2025)]{drayson2025contamination}
Ryan Drayson, Laura Forti, and others.
\newblock Preventing model collapse via importance resampling of synthetic data.
\newblock \emph{arXiv preprint}, 2025.

\bibitem[Mitchell et~al.(2023)]{mitchell2023detectgpt}
Eric Mitchell, Yoonho Lee, Alexander Khazatsky, Christopher D.~Manning, and Chelsea Finn.
\newblock {DetectGPT}: Zero-shot machine-generated text detection using probability curvature.
\newblock \emph{ICML}, 2023.

\bibitem[Xie et~al.(2023)]{xie2023dsir}
Sang~Michael Xie, Shibani Santurkar, Tengyu Ma, and Percy Liang.
\newblock Data selection for language models via importance resampling.
\newblock \emph{NeurIPS}, 2023.

\end{thebibliography}

\end{document}
